\chapter{Gaussian Quadrature}
\label{ch:17}
\noindent\emph{Source and derivation note.}
The orthogonal system and normalization are those of Sections~\ref{sec:1-2}, \ref{sec:5-2}, and \ref{sec:13-1}. The product-martingale source is the supplied \texttt{sources/KravFeinsKocik.tex}, identified in the source register. The Gaussian rule, terminal factorization, proofs, and numerical examples here are developed from those identities; they are not attributed to an unread item in the supplied literature list. Numerical provenance and reproduction details are in \texttt{sources/INTEGRATION\_gaussian\_summation.md} and \texttt{sources/compute\_gaussian\_krawtchouk.py}.

\section{Zeros of Krawtchouk polynomials}
\label{sec:17-1}
The measure in this chapter is finite and discrete:
\[
\mu_N=2^{-N}\sum_{j=0}^N\binom Nj\delta_j,\qquad
\mathcal S_N[f]=2^{-N}\sum_{j=0}^N\binom Nj f(j).
\]
Thus Gaussian quadrature here compresses a weighted sum, not an ordinary integral. Keep two integers distinct: $N$ specifies the $N+1$ support points, whereas $n$ specifies the degree of the orthogonal polynomial and the number of Gaussian nodes. We take $1\le n\le N$ unless stated otherwise.

Extend $K_n(j;N)$ polynomially off the integer lattice via the recurrence of Section~\ref{sec:1-3}. Its leading coefficient in $j$ is $(-2)^n/n!$. The monic polynomials
\[
p_n(j)=\frac{(-1)^n n!}{2^n}K_n(j;N)
=\frac{(-1)^n}{2^n}k_n(N-2j,N)
\]
therefore satisfy
\begin{equation}\label{eq:quadrature-monic-recurrence}
p_{r+1}(j)=(j-N/2)p_r(j)-\frac{r(N-r+1)}4p_{r-1}(j),\qquad p_0=1.
\end{equation}
Their squared norms under $\mu_N$ are
\begin{equation}\label{eq:quadrature-monic-norm}
h_r=\mathcal S_N[p_r^2]=\frac{(r!)^2}{4^r}\binom Nr
=\frac{r!N!}{4^r(N-r)!}.
\end{equation}
The Gaussian nodes $z_1<\cdots<z_n$ are the zeros of $p_n$, equivalently the eigenvalues of the $n\times n$ real symmetric Jacobi matrix
\begin{equation}\label{eq:quadrature-Jacobi}
\mathcal J_{N,n}=\begin{pmatrix}
N/2&b_1& & &\\
b_1&N/2&b_2& &\\
 &b_2&N/2&\ddots&\\
 & &\ddots&\ddots&b_{n-1}\\
 & & &b_{n-1}&N/2
\end{pmatrix},\qquad b_r=\frac12\sqrt{r(N-r+1)}.
\end{equation}
For $n=1$ this means the scalar matrix $(N/2)$. Positivity of $b_r$ gives simple real eigenvalues; the Rayleigh quotient is the weighted average of $j$ against the square of a nonzero polynomial of degree at most $n-1$. Such a polynomial cannot vanish on all $N$ other support points, so $0<z_a<N$. The zeros of successive degrees strictly interlace. Reflection symmetry gives
$z_{n+1-a}=N-z_a$.
These nodes are generally not integers: polynomial extension, rather than a table of lattice values alone, is essential.

\section{Integral zeros}
\label{sec:17-2}
\paragraph{The trivial terminal zeros.}
Remark~\ref{rem:terminal-integral-zeros} gives a universal family of integral zeros at the terminal degree:
\[
k_{N+1}(x,N)=\prod_{j=0}^{N}(x-(N-2j)).
\]
For even $N=2m$ this is $x\prod_{r=1}^{m}(x^2-4r^2)$; for odd $N=2m+1$ it is $\prod_{r=0}^{m}(x^2-(2r+1)^2)$. Thus all centered spectral points are zeros. For example,
\[
k_6(x,5)=(x^2-1)(x^2-9)(x^2-25).
\]
In the uncentered coordinate $j=(N-x)/2$, this is precisely the monic terminal polynomial $p_{N+1}(j)=\prod_{r=0}^{N}(j-r)$ discussed below. This terminal case is automatic and should be distinguished from the nontrivial question of integral zeros in degrees $1\le n\le N$.

Integer zeros are special cases of these polynomial zeros, not a prerequisite for Gaussian summation. Since
$K_n(N-j;N)=(-1)^nK_n(j;N)$, every odd degree has the central zero $j=N/2$; this is integral exactly when $N$ is even. For degree two,
\[
K_2(j;N)=\frac{(N-2j)^2-N}{2},\qquad
z_\pm=\frac{N\pm\sqrt N}{2}.
\]
For $N\ge2$ these two zeros are integers exactly when $N$ is a perfect square: a square and its square root have the same parity. For example, $N=9$ gives nodes $3$ and $6$, each with Gaussian weight $1/2$. Consequently the two-node rule integrates, in the discrete sense, every cubic exactly under $\mu_9$.
For $N=5$ the degree-three centered polynomial is $k_3(x,5)=x(x^2-13)$, giving $j=(5-\sqrt{13})/2,5/2,(5+\sqrt{13})/2$, none integral.
These examples are not a classification of integral zeros at higher degree.
\par\noindent\textbf{A GENERAL CLASSIFICATION OF INTEGRAL ZEROS NEEDS TO BE DONE.}

\section{Gaussian--Krawtchouk summation}
\label{sec:17-3}
For each Gaussian node define the Lagrange polynomial
\[
\ell_a(j)=\frac{p_n(j)}{(j-z_a)p_n'(z_a)},\qquad
\lambda_a=\mathcal S_N[\ell_a].
\]
Then the $n$-node formula is
\begin{equation}\label{eq:gaussian-Krawtchouk-sum}
2^{-N}\sum_{j=0}^N\binom Nj f(j)
=\sum_{a=1}^n\lambda_a f(z_a),\qquad \deg f\le2n-1.
\end{equation}
To prove exactness, divide $f=q p_n+r$ with $\deg q,\deg r\le n-1$. Orthogonality annihilates $\mathcal S_N[q p_n]$, while interpolation of $r$ evaluates its sum exactly. Since $p_n(z_a)=0$, the right side is also unchanged by removing $q p_n$.
Applying exactness to $\ell_a^2$ shows
$\lambda_a=\mathcal S_N[\ell_a^2]>0$, and applying it to $1$ gives $\sum_a\lambda_a=1$. The rule is symmetric: $\lambda_{n+1-a}=\lambda_a$.
Two equivalent computable weight formulas are
\begin{equation}\label{eq:quadrature-Christoffel}
\lambda_a=\left[\sum_{r=0}^{n-1}\frac{K_r(z_a;N)^2}{\binom Nr}\right]^{-1}
=\frac{h_{n-1}}{p_n'(z_a)p_{n-1}(z_a)}.
\end{equation}
Alternatively, if $u_a$ is the unit eigenvector of \eqref{eq:quadrature-Jacobi} for $z_a$, then $\lambda_a=|(u_a)_0|^2$. The spectral measure of the first coordinate gives this identity; the Christoffel formula follows by summing squares of orthonormal polynomials. These formulas use the probability-normalized binomial measure. For the unnormalized sum $\sum_j\binom Nj f(j)$ multiply every weight by $2^N$.

\subsection{Error and the finite-support endpoint}
For $f\in C^{2n}([0,N])$, Hermite interpolation at the $n$ nodes, followed by summation against the positive measure, gives
\[
\mathcal S_N[f]-\sum_a\lambda_a f(z_a)
=\frac{h_n}{(2n)!}f^{(2n)}(\xi),\qquad \xi\in[0,N].
\]
In particular, the error on the monic polynomial $j^{2n}$ is exactly $h_n>0$, so degree $2n-1$ is sharp. The rule remains meaningful for nonpolynomial functions provided values between lattice points are defined. If only the $N+1$ discrete data values are known, choosing an off-lattice interpolation is extra information; one cannot assume an arbitrary data array has prescribed Gaussian-node values.
For $n=N+1$, the terminal polynomial is $p_{N+1}(j)=\prod_{r=0}^N(j-r)$ and its norm is zero. The full Jacobi matrix recovers the original nodes $0,\ldots,N$ and weights $2^{-N}\binom Nj$. This is the original finite sum, exact for every lattice function; it is not a new compressed rule, and the positive-norm error discussion above is restricted to $n\le N$.

\subsection{Accurate examples for $N=5,10,15$}
The following tables give both reduced-node rules $(N,n)=(5,3),(10,5),(15,8)$ and the degree-$N$ rules $(5,5),(10,10),(15,15)$. Thus the reduced rules replace respectively $6,11,16$ summands by $3,5,8$ evaluations and are exact through degrees $5,9,15$. The degree-$N$ rules are exact through degrees $9,19,29$. All nodes below are in the original $j$ coordinate; centered nodes are $x_a=N-2z_a$ with the same weights.

For $(N,n)=(5,3)$ one also has the exact formula
\[
\mathcal S_5[f]=\frac5{26}f\!\left(\frac{5-\sqrt{13}}2\right)
+\frac8{13}f\!\left(\frac52\right)
+\frac5{26}f\!\left(\frac{5+\sqrt{13}}2\right),\qquad\deg f\le5.
\]
The tables show 18 significant digits. Their machine-readable companions retain 65 digits. Computation used a symmetric Jacobi eigensystem at 90-decimal working precision, independently checked against the centered polynomial recurrence, the Christoffel weight formula, positivity, and every scaled moment $(j/N)^r$ for $0\le r\le2n-1$. Maximum absolute scaled-moment error over all six rules was below $10^{-85}$ at working precision. These numerical residuals are consistency checks, not interval-certified error bounds; the printed rounded tables have correspondingly lower accuracy.
\par\medskip\noindent\textbf{$N=5$, $n=3$; exact through degree $5$.}
\begin{center}\small
\begin{tabular}{rll}
$a$ & $z_a$ & $\lambda_a$\\\hline
1 & 0.697224362268005327 & 0.192307692307692318\\
2 & 2.5 & 0.615384615384615419\\
3 & 4.30277563773199478 & 0.192307692307692318\\
\end{tabular}
\end{center}
\par\medskip\noindent\textbf{$N=5$, $n=5$; exact through degree $9$.}
\begin{center}\small
\begin{tabular}{rll}
$a$ & $z_a$ & $\lambda_a$\\\hline
1 & 0.0649539076702599777 & 0.0406121317521831102\\
2 & 1.24678392596102428 & 0.244622767576675948\\
3 & 2.5 & 0.42953020134228187\\
4 & 3.7532160740389755 & 0.244622767576675948\\
5 & 4.93504609232974012 & 0.0406121317521831102\\
\end{tabular}
\end{center}
\par\medskip\noindent\textbf{$N=10$, $n=5$; exact through degree $9$.}
\begin{center}\small
\begin{tabular}{rll}
$a$ & $z_a$ & $\lambda_a$\\\hline
1 & 1.0 & 0.01953125\\
2 & 3.0 & 0.234375\\
3 & 5.0 & 0.4921875\\
4 & 7.0 & 0.234375\\
5 & 9.0 & 0.01953125\\
\end{tabular}
\end{center}
\par\medskip\noindent\textbf{$N=10$, $n=10$; exact through degree $19$.}
\begin{center}\small
\begin{tabular}{rll}
$a$ & $z_a$ & $\lambda_a$\\\hline
1 & 0.00426140162761967471 & 0.00100005627585995687\\
2 & 1.03215639227186573 & 0.0108643711646248114\\
3 & 2.1094239579268903 & 0.0553578832768971893\\
4 & 3.24085911987510711 & 0.160428693955712715\\
5 & 4.40940159989234548 & 0.272348995326905319\\
6 & 5.59059840010765452 & 0.272348995326905319\\
7 & 6.75914088012489334 & 0.160428693955712715\\
8 & 7.8905760420731097 & 0.0553578832768971893\\
9 & 8.96784360772813471 & 0.0108643711646248114\\
10 & 9.99573859837238032 & 0.00100005627585995687\\
\end{tabular}
\end{center}
\par\medskip\noindent\textbf{$N=15$, $n=8$; exact through degree $15$.}
\begin{center}\small
\begin{tabular}{rll}
$a$ & $z_a$ & $\lambda_a$\\\hline
1 & 0.810168654648509734 & 0.000551841058012324075\\
2 & 2.70329267453015554 & 0.0178684857525261559\\
3 & 4.61682155336053412 & 0.134507504747883117\\
4 & 6.53821132811295858 & 0.347072168441578399\\
5 & 8.46178867188704054 & 0.347072168441578399\\
6 & 10.3831784466394659 & 0.134507504747883117\\
7 & 12.2967073254698445 & 0.0178684857525261559\\
8 & 14.1898313453514895 & 0.000551841058012324075\\
\end{tabular}
\end{center}
\par\medskip\noindent\textbf{$N=15$, $n=15$; exact through degree $29$.}
\begin{center}\small
\begin{tabular}{rll}
$a$ & $z_a$ & $\lambda_a$\\\hline
1 & 0.000210675294063239482 & 0.0000305592264942427708\\
2 & 1.00264068454641486 & 0.000462988590523494585\\
3 & 2.01486605080593995 & 0.00335579980348639896\\
4 & 3.05052871756057487 & 0.0155659622211499205\\
5 & 4.12061368475752943 & 0.0506438220470849867\\
6 & 5.22535900511131501 & 0.117850851005629698\\
7 & 6.35599416357050373 & 0.195969221661716791\\
8 & 7.5 & 0.232241590887828919\\
9 & 8.64400583642949627 & 0.195969221661716791\\
10 & 9.77464099488868499 & 0.117850851005629698\\
11 & 10.8793863152424706 & 0.0506438220470849867\\
12 & 11.9494712824394256 & 0.0155659622211499205\\
13 & 12.9851339491940596 & 0.00335579980348639896\\
14 & 13.9973593154535845 & 0.000462988590523494585\\
15 & 14.9997893247059366 & 0.0000305592264942427708\\
\end{tabular}
\end{center}

